% Vorlage für WABs der Provadis Hochschule
% basiert auf "Universität Ulm Praktikumsbericht Vorlage"
% von Max Sch.
% CC BY 4.0
%
% 02_Stand.tex

\chapter{Stand der Forschung}

In diesem Kapitel wird besonders intensiv mit der Literatur gearbeitet. Entsprechend werden hier sehr viele Zitate vorkommen. Ich werde die Quellenhinweise sofort einbauen, damit ich sie nicht vergesse.

Dabei beachte ich die Vorgaben der Provadis School of International Management and Technology, die in den „Anforderungen an wissenschaftliche Hausarbeiten“ beschrieben sind.

Die Literatur wird im Literaturverzeichnis am Ende der Arbeit aufgeführt. Es enthält alle verwendeten Quellen -- und nur diese --, auch Internetquellen (mit Webadresse in eckigen Klammern und Datum des letzten Zugriffs).

Literaturquellen werden in der Datei Referenzen.bib im BibTeX-Format eingetragen und zitiert. Die Seitenangabe ist obligatorisch.
\ifharvard
    Ein wörtliches Zitat steht in Anführungszeichen: \enquote{Kaffee ist die einzige Substanz, die uns schneller arbeiten lässt}\footcite[12]{lustigeCitation}. Eine Paraphrase steht ohne Anführungszeichen, die Fußnote wird mit \enquote{Vgl.} eingeleitet\footcite[Vgl.][8\psq]{mustermann2020}. Auch Internetquellen werden so belegt\footcite[Vgl.][]{provadisWeb}.
\else
    Ein wörtliches Zitat steht in Anführungszeichen: \enquote{Kaffee ist die einzige Substanz, die uns schneller arbeiten lässt}~\cite[12]{lustigeCitation}. Paraphrasen stehen ohne Anführungszeichen~\cite[8\psq]{mustermann2020}.
\fi

Wörtlich übernommene Stellen aus KI-Quellen stehen in Asterisken: \kizitat{Ein gutes Recruiting beginnt mit einer klaren Stellenbeschreibung.}{ChatGPT 1} Wurde KI-Output nur als Ausgangspunkt verwendet, lautet die Fußnote \enquote{Vgl. Elicit 1.}\kivgl{Elicit 1}

Gut gepflegte BibTeX-Einträge finden sich etwa in der DBLP\footnote{\url{https://dblp.uni-trier.de}} und in den Online-Bibliotheken großer Wissenschaftsverlage, wie der ACM Digital Library\footnote{\url{https://dl.acm.org}}, IEEE Xplore\footnote{\url{https://ieeexplore.ieee.org}}, Springer Link\footnote{\url{https://link.springer.com}}, Elsevier ScienceDirect\footnote{\url{https://www.elsevier.com/de-de/products/sciencedirect}}, uvw.
\todo[inline]{Quellenverweise sollten immer \emph{direkt} hinter Nennung der Quelle bzw. der zu unterstützenden Aussage stehen, und nicht am Ende des Satzes oder gar des Absatzes.}

Weitere Verzeichnisse sind alphabetisch zu sortieren und besonders sorgfältig anzufertigen.
